%%%%%%%%%%%%%%%%%%%%%%%%%%%%%%%%%%%%%%%%%%%%%
\section{Conclusions}
\label{Sec:Conclusions}
%%%%%%%%%%%%%%%%%%%%%%%%%%%%%%%%%%%%%%%%%%%%%

In these notes we have studied numerically the relaxation through phase space mixing of a spherically symmetric collisionless gas which is trapped in an external central potential, taking into account its self-gravity, in the reduced model in which all the gas particles have the same magnitude of the angular momentum. Our analysis is based on a PIC code whose particles are placed on a regular lattice in action-angle variables, on a solution of the linearized equations in the action-angle variables of a self-consistent steady state, and on the gain $\lambda(\omega)$ of the self-consistent loop. In this section we summarize the results, point out a feature which is common to them, and list some open problems.

For initial data close to a self-consistent steady state, the self-gravity of the deviation from that state has two effects. Below the threshold mass at which $\lambda_{edge} = 1$, the deviation decays, first like a damped oscillation and later algebraically, but more slowly than in a fixed potential, and at late times, or for larger masses, the exponent of the algebraic tail is reduced from $k + 1$ to $k$. Above the threshold the deviation contains an oscillating mode, whose frequency agrees with the root of $\lambda(\omega_d) = 1$ to within $2\times 10^{-4}$ in relative terms, such that the mixing is not completed. At finite amplitude, the gas particles near the edge of the steady state are trapped by the mode when the amplitude exceeds the critical value~(\ref{Eq:epsc}), which vanishes quadratically with the distance of the mode from the band, and the trapping reduces the amplitude of the mode; below the threshold, the same mechanism stops the damping of the deviation, as predicted by O'Neil for plasma waves~\cite{tO65}.

For initial data far from equilibrium, in which the gas is concentrated about one orbital phase, the self-gravity of the gas first slows down the mixing and then stops it. When the parameter~(\ref{Eq:muDef}), which compares the bounce frequency in the potential well of the cloud with the spread of its orbital frequencies, is smaller than a critical value $\mu_c$, the gas mixes and its potential becomes stationary, whereas above $\mu_c$ it collapses into a compact, self-bound clump which oscillates radially about the circular orbit and does not mix during the $30$ units of the mixing time $\tau_1$ covered by our simulations. For clouds which are compact in angle, $\mu_c$ lies between $0.43$ and $0.48$ in the isochrone potential as well as in the Kepler potential, independently of the width of the cloud in action. In the Kepler potential this means that, for a cloud of a given shape, phase mixing in the sense of Ref.~\cite{sCjL26} only takes place below a critical mass.

The two situations share a common feature. In both of them the outcome of the competition between mixing and self-gravity is controlled by the ratio between a bounce frequency, which measures the strength of the potential that holds a group of particles together, and a difference of orbital frequencies, which measures the shear that separates them: the ratio~(\ref{Eq:sDef}) between the bounce frequency in the potential of the mode and the distance of the mode from the band, and the parameter $\mu$ for the cloud. In both cases a simple pendulum model gives the correct order of magnitude of the critical value, and in the case of the mode near the edge it also gives the largest action variable which is reached by the trapped particles, without adjustable parameters.

The numerical methods required some care. The functions $h_n(t)$ have to be computed with the action-angle variables of the total potential, or of the instantaneous potential for data far from equilibrium, since the action-angle variables of the external potential leave a time-independent residue which could be mistaken for an incomplete mixing. The quiet start on a lattice in action-angle variables reduces the noise by several orders of magnitude, but its regularity is lost in the course of time due to the discreteness of the particles; for $M_{gas}/M = 1$ about $10^5$ particles and amplitudes $\varepsilon\leq 0.01$ are needed in order to reproduce the solutions of the linearized equations within a few percent, and close to the transition of the cloud the results depend on the number of particles.

Several questions remain open. For the data close to a steady state, the derivation of the tail $t^{-k}$ is formal, and its amplitude is only known to leading order in $M_{gas}/M$; for the modes near the edge, the depth of the loss of amplitude caused by the trapped particles is not explained by the pendulum model. For the cloud, a theory of the critical value $\mu_c$ is missing, which should also explain why a cloud which is wide in angle requires a larger value; the mass of the clump has to be defined by a dynamical criterion; and it is not known whether the oscillation of the clump is eventually damped, as in the core-halo states of systems with long-range interactions~\cite{yLrPfRtTfB14}, and on which time scale. Answering the last question requires longer simulations with more particles. Finally, all our results refer to the reduced model. When the angular momenta of the gas particles are distributed, the frequencies of the orbits depend also on $L$, which provides an additional source of mixing (see section~\ref{SubSec:FreeMixing}); whether this shifts the critical values found here is a natural next step, for which the methods of these notes can be extended in a direct way. \EJ{To be completed after the discussion with Olivier: the outlook regarding the distribution of angular momenta, and whether the material of section~\ref{Sec:FixedL} on the threshold mass is kept here or moved to a separate article.}
